\documentclass{ieeeaccess}
\usepackage{cite}
\usepackage{amsmath,amssymb,amsfonts}
\usepackage{algorithmic}
\usepackage{graphicx}
\usepackage{textcomp}
\usepackage{xcolor}
\usepackage{array}
\interdisplaylinepenalty=2500

\begin{document}

\history{Date of publication October 29, 2026, date of current version October 29, 2026.}
\doi{10.1109/ACCESS.2026.XXXXXXX}
\vol{XX}
\year{2026}

\title{Rethinking the Reported Accuracy of Smartphone-Based Electric-Field Inversion for Crowdsourced EMF Monitoring}

\author{\uppercase{Lianzhong Zheng}\authorrefmark{1},
\uppercase{and Juan Pan}\authorrefmark{1}}

\address[1]{China Academy of Information and Communications Technology, Beijing 100191, China (e-mail: zlz347661070@gmail.com, panjuan@caict.ac.cn).}

\tfootnote{This work was not supported by any external funding. The authors declare that they have no known competing financial interests or personal relationships that could have appeared to influence the work reported in this paper.}

\markboth
{Zheng \headeretal: Rethinking the Reported Accuracy of Smartphone-Based Electric-Field Inversion}
{Zheng \headeretal: Rethinking the Reported Accuracy of Smartphone-Based Electric-Field Inversion}

\corresp{Corresponding author: Juan Pan (e-mail: panjuan@caict.ac.cn).}

\begin{abstract}
Smartphone-based inversion of electric field strength is attractive for
crowdsourced electromagnetic environment monitoring, yet reported mean absolute
errors (MAE) around 1~dB are almost always obtained with a random train/test split
of geo-referenced 1~Hz traces, a protocol in which neighbours of every test sample
remain in the training set. Using a public dataset of 46{,}306 samples, we
reproduce a representative gradient-boosting study and then audit how accuracy
depends on three rarely varied factors. First, \emph{protocol}: the same model
yields 0.97, 1.05 and 1.35~dB under a random hold-out, spatio-temporal patch-based
cross-validation with an embargo, and leave-one-group-out, respectively;
protocol choice moves the estimate by 0.39~dB, more than twice the spread among
six regressors. Second, \emph{features}: adding the reported EARFCN (E-UTRA absolute radio frequency channel number)
reduces MAE from 1.05 to 0.73~dB and raises the weakest-band $R^2$ from $-0.41$
to 0.65; the size of the gain depends on how the label is defined and ranges
from 15\% to 31\% across three definitions.
Third, \emph{device calibration}: input-space recalibration of the same handset
is absorbed by the tree ensemble, whereas a 2~dB inter-device offset degrades MAE
by 30\% and a 100-sample output calibration recovers 65--75\% of that loss.
Even under the strictest protocol the band-aware model still outperforms the
originally reported random-split result, by 10\% (25\% under spatio-temporal
blocking). We conclude that protocol, feature set and
device identity must be reported together, and give recommendations for
crowdsourced deployment.
\end{abstract}

\begin{keywords}
Electromagnetic field exposure, crowdsourced monitoring, machine learning,
cross-validation, data leakage, LTE, reference signal received power.
\end{keywords}

\titlepgskip=-21pt

\maketitle

\section{Introduction}
\label{sec:introduction}
\PARstart{C}{rowdsourced} monitoring of radio-frequency electromagnetic fields
(EMF) promises a coverage and a cost that fixed monitoring stations cannot match\cite{aerts2013,ganti2011,lane2010}.
Regulatory frameworks such as the ICNIRP guidelines \cite{icnirp} and IEEE C95.1\cite{ieee2019} define exposure
limits in terms of the electric field strength $E$, which a smartphone cannot
measure directly: the handset reports downlink quality indicators, most notably
the reference signal received power (RSRP), and a calibration model must map them
to $E$.

Early approaches used a fixed conversion factor derived from standardized
extrapolation procedures \cite{schie} and were limited to roughly 5~dB of
accuracy in complex propagation environments. Subsequent work replaced the fixed
factor with a learned mapping. Zeng \emph{et al.} \cite{zeng2024} trained a
gradient boosting decision tree (GBDT) \cite{friedman,friedman2002,hastie2009} on 23 features derived
from a seven-sample RSRP window and reported a coefficient of determination
$R^2=0.93$ and MAE~$=0.97$~dB, an improvement of about 4~dB over the fixed
factor.
More recently, Huang \emph{et al.} \cite{huang2026} fused RSRP, RSSI, RSRQ and
SINR in a stacking ensemble and reached 0.71~dB, attributing the gain to the
complementary information carried by the additional parameters.

Both results were obtained with a random 80/20 split of data acquired by walking a
route at 1~Hz. Consecutive samples are then about 1.3~m and 1~s apart, so the
neighbours of a test sample almost certainly remain in the training set. This
inflates accuracy in exactly the regime that a crowdsourcing deployment must
handle, namely predicting at locations and times that were never measured.
Spatial and temporal blocking are standard practice in geostatistics and remote
sensing \cite{roberts,kara,pohjankukka2017,meyer2021,valavi2019}, but have not, to our knowledge, been applied to
smartphone-based EMF inversion, and the magnitude of the resulting bias is
unknown.

A second observation motivates this work. The handset knows which carrier
frequency it is receiving: the E-UTRA absolute radio frequency channel number
(EARFCN) is reported by the modem and is available at inference time. The baseline
model ignores it, although the four measured bands differ by 12.4~dB in mean field
strength and are even labelled with different statistics. Exploiting a quantity
that is already available is free, and we quantify what it is worth.

This paper makes four contributions.
\begin{enumerate}
\item We reproduce a representative single-parameter EMF inversion model on its
public dataset to within 0.01 in $R^2$ and 0.02~dB in MAE, and recover
implementation details that the original paper does not report
(Appendix~\ref{app:details}).
\item We quantify how much of the reported accuracy is an artefact of the
evaluation protocol, using three protocols of increasing strictness, and show that
protocol choice dominates model choice. We further show that $R^2$ is nearly
insensitive to leakage, and that per-fold $R^2$ becomes meaningless when a fold contains a single band, because the internal variance is too small.
\item We show that a single EARFCN feature recovers most of the accuracy that a
band-agnostic model loses, and that a per-band model gives no further gain.
\item We show that input-space recalibration of the same device is absorbed by the
model, whereas cross-device calibration is essential, and we quantify how many
labelled samples it requires.
\end{enumerate}

The remainder of the paper is organised as follows. Section~\ref{sec:data}
describes the dataset and the reproduction. Section~\ref{sec:protocol} defines the
three protocols. Section~\ref{sec:results} reports the results and ranks the
factors. Section~\ref{sec:device} treats device calibration.
Section~\ref{sec:discussion} discusses implications and
Section~\ref{sec:conclusion} concludes.

\section{Dataset and Baseline Reproduction}
\label{sec:data}

\subsection{Measurement Setup}
The public dataset of \cite{zeng2024} was acquired on a campus area of about
$380\times260$~m during five campaigns on three days in May and June 2023, with
four Long Term Evolution (LTE) downlink bands (EARFCN 100, 1850, 1300 and 38400, i.e. 2110--2130,
1860--1880, 1805--1825 and 1885--1905~MHz). A smartphone application logged RSRP
(defined in 3GPP TS~36.214 \cite{3gpp}) at approximately 1~Hz while a spectrum analyzer with an isotropic antenna scanned
the corresponding downlink band; the two streams were aligned by timestamp. The
released sample table contains 46{,}306 samples, and every sample carries a GPS
fix (100\% match), which is what makes the blocking schemes of
Section~\ref{sec:protocol} possible.

\subsection{Feature Construction}
For each field-strength timestamp, the RSRP series is cleaned at $-110$~dBm,
smoothed by an exponentially weighted moving average (EWMA) and resampled on a 1~s
grid with linear interpolation of gaps shorter than 5~s. A seven-point window
($t-3,\ldots,t+3$) then yields 16 statistics---mean, minimum, maximum, median,
first and third quartiles, range, variance, standard deviation, mode, skewness,
kurtosis, mean and median absolute deviation, interquartile range and coefficient
of variation---to which the seven raw RSRP values are appended, giving 23
features. The label is the field strength measured by the spectrum analyzer, taken
as the third quartile for EARFCN~100 and as the maximum for the other three bands.

\subsection{Reproduction}
We trained the published configuration (Huber loss with $\delta=0.9$, learning
rate 0.05, maximum depth 5, 200 trees, 80/20 split) on the released sample table.
Table~\ref{tab:repro} shows that all four metrics match the published values to
within 0.02~dB and 0.01 in $R^2$.

\begin{table}
\caption{Reproduction of the baseline model (random 80/20 split).}
\label{tab:repro}
\setlength{\tabcolsep}{3pt}
\begin{tabular}{|l|c|c|c|c|}
\hline
 & $R^2$ & MSE & RMSE (dB) & MAE (dB)\\
\hline
Training set, reported \cite{zeng2024} & 0.94 & 1.48 & 1.22 & 0.93\\
Training set, reproduced & 0.9387 & 1.4773 & 1.2154 & 0.9312\\
Test set, reported \cite{zeng2024} & 0.93 & 1.60 & 1.26 & 0.97\\
Test set, reproduced & 0.9346 & 1.5657 & 1.2513 & 0.9574\\
\hline
\end{tabular}
\end{table}

Rebuilding the whole feature table from the raw measurements reproduces the
released seven raw RSRP values exactly (MAE~$=0$) and the positional statistics
with a correlation of 0.988, which confirms that the pipeline of
Section~\ref{sec:data}-B is the one that generated the released samples. Details
that had to be inferred are listed in Appendix~\ref{app:details}.

\section{Evaluation Protocols}
\label{sec:protocol}
All protocols use the same model and the same features; only the partition
changes. Fig.~\ref{fig:map} shows the study area and the two block-wise
partitions.

\Figure[t!](topskip=0pt, botskip=0pt, midskip=0pt){figures/fig3_study_area.pdf}
{\textbf{Study area with samples coloured by fold (five folds).} (a) 50~m spatial
patches: a whole patch belongs to one fold. (b) Patches crossed with 10~min time
windows: the same patch appears in several folds at different times.\label{fig:map}}

\subsection{Random Hold-out (RH)}
The baseline protocol: 20\% of the samples drawn uniformly at random, repeated
over several seeds. Test samples retain temporal neighbours (1~s) and spatial
neighbours (about 1.3~m) in the training set. Measured on the reproduction split,
99.9\% of the test samples lie within 2~m of a training sample and the median
time gap to the nearest training sample is 1~s.

\subsection{Spatio-Temporal Patch-Based CV with Embargo (ST-PB-CV)}
The area is divided into 50~m patches and the timeline into 10~min blocks; the
group is the (patch, block) pair, and groups are assigned to five folds. An
embargo removes from the training set any sample lying within 20~m \emph{and}
60~s of a test sample, suppressing residual autocorrelation across block
boundaries. The block size is matched to the physics of the features: 50~m is
about 40~s of walking, comparable to the 3~s window and the $\approx$20~s EWMA
time constant of the $\alpha=0.05$ filter.

\subsection{Leave-One-Group-Out CV (LOGO-CV)}
The group is the (campaign, band) pair, giving ten folds. Holding out one group
removes an entire measurement session of one band, which mimics onboarding a new
campaign without recalibration. Grouping by campaign alone would be incorrect
here: each campaign measured only two of the four bands, so a campaign-level
hold-out would simultaneously remove two bands from training and confound
cross-session with cross-band generalisation. We verified that the campaign-level
variant yields 1.33~dB, comparable to the (campaign, band) variant (1.35~dB) reported
below.

\subsection{Metrics}
We report MAE, MSE (mean squared error) and RMSE (root mean squared error), and a \emph{pooled} $R^2$ computed once over the
concatenated test predictions of all folds. Per-fold $R^2$ is not comparable
across folds when the folds have different label variance, as shown in
Section~\ref{sec:r2}.

\section{Results}
\label{sec:results}

\subsection{Protocol Effect}
Table~\ref{tab:proto} and Fig.~\ref{fig:proto} summarise the three protocols. MAE
increases from 0.966~dB (RH) to 1.046~dB (ST-PB-CV) and 1.353~dB (LOGO-CV). The
gap between the loosest and the strictest protocol is 0.39~dB, or 40\%.

\begin{table}
\caption{Effect of the evaluation protocol on the same model and features.}
\label{tab:proto}
\setlength{\tabcolsep}{3pt}
\begin{tabular}{|l|c|c|c|c|}
\hline
Protocol & Folds & MAE (dB) & RMSE (dB) & pooled $R^2$\\
\hline
RH (8 seeds) & 8 & 0.966 & 1.261 & 0.934\\
ST-PB-CV + embargo & 5 & 1.046 & 1.351 & 0.924\\
LOGO-CV & 10 & 1.353 & 1.637 & 0.886\\
\hline
\end{tabular}
\end{table}

\Figure[t!](topskip=0pt, botskip=0pt, midskip=0pt){figures/fig1_protocols.pdf}
{\textbf{Test MAE under the three protocols.} Error bars are 95\% confidence
intervals over folds or seeds.\label{fig:proto}}

\Figure[t!](topskip=0pt, botskip=0pt, midskip=0pt){figures/fig_pred_scatter.pdf}
{\textbf{Predicted against measured field strength.} (a) Random hold-out, using
the reproduction split of Table~\ref{tab:repro}. (b) LOGO-CV, i.e., leaving one
(campaign, band) group out. The
dashed line is $y=x$; open squares mark the mean of each band. The cloud widens
and the band means drift away from the identity line as the protocol is tightened.\label{fig:scatter}}

Fig.~\ref{fig:scatter} shows the same effect sample by sample. Under the random
hold-out the cloud of predictions is tight around the identity line and the
per-band means stay close to it. Under LOGO-CV the cloud widens and the band
means drift off the line, reaching $+0.98$~dB for EARFCN~100, while the global
bias stays small ($+0.11$~dB). The loss is therefore dispersion plus a
band-dependent shift, not a constant offset that a single global correction could
remove.

\subsection{Why the Coefficient of Determination Misleads}
\label{sec:r2}
Pooled $R^2$ barely changes across protocols (0.934, 0.924 and 0.886) while MAE
changes by 40\%. The reason is that $R^2$ is dominated by the between-band spread
of the field strength, which blocking does not remove. Fig.~\ref{fig:cdf} shows
the two mechanisms at work. Under the random split the training and test CDFs of
the field strength are indistinguishable (Kolmogorov--Smirnov $D=0.012$,
$p=0.27$): the marginals match exactly, so a distributional check cannot detect
the leakage, which resides in the joint spatio-temporal structure. Under LOGO-CV
the ten (campaign, band) groups each occupy a narrow slice of the 13.3~dB spanned
by their means, so a held-out group is internally quiet
($\sigma=0.73$--$2.81$~dB). An error of 1.4~dB then produces a negative per-fold
$R^2$: the values range from $-2.62$ to $0.61$ with a mean of $-0.18$, even
though the MAE is unremarkable. Pooling the predictions restores a meaningful
value (0.886). We therefore recommend that this literature report MAE and RMSE,
and pooled $R^2$ if a coefficient of determination is required.

\Figure[t!](topskip=0pt, botskip=0pt, midskip=0pt){figures/fig_field_cdf.pdf}
{\textbf{Marginal distribution of the field strength.} (a) Under a random 80/20
split the training and test CDFs coincide (two-sample Kolmogorov--Smirnov
$D=0.012$, $p=0.27$); distributional similarity is not evidence that the split is
sound. (b) The ten (campaign, band) groups used by LOGO-CV each occupy a narrow
slice of the range, which is why per-fold $R^2$ is unstable.\label{fig:cdf}}

\subsection{Band Information}
The baseline model is band-agnostic, and this, rather than the information content
of RSRP, explains why its within-band accuracy is poor. Adding EARFCN as a single
numeric feature reduces MAE from 1.046 to 0.730~dB under ST-PB-CV
(Table~\ref{tab:band}, Fig.~\ref{fig:ours}) and raises the pooled $R^2$ from 0.924
to 0.958. Training four independent models, one per band and selected at inference
by the current EARFCN, gives 0.737~dB, a difference well inside the fold-to-fold
spread ($\pm0.04$~dB). Since a tree ensemble can already split on EARFCN, a
per-band model duplicates information that a single model obtains for free, at
four times the maintenance cost.

\begin{table}
\caption{Protocol $\times$ feature design. Entries are test MAE in dB.}
\label{tab:band}
\setlength{\tabcolsep}{3pt}
\begin{tabular}{|l|c|c|c|}
\hline
 & RH & ST-PB-CV & LOGO-CV\\
\hline
RSRP only (baseline) & 0.969 & 1.046 & 1.353\\
$+$ EARFCN (single model) & 0.667 & 0.730 & 0.871\\
Per-band models (four) & 0.628 & 0.737 & 0.868\\
\hline
\end{tabular}
\end{table}

\Figure[t!](topskip=0pt, botskip=0pt, midskip=0pt){figures/fig2_protocol_feature.pdf}
{\textbf{Two independent factors.} Even under the strict protocol, the band-aware
model (0.730~dB) beats the originally reported random-split value (0.97~dB).\label{fig:ours}}

The within-band $R^2$ changes accordingly (Table~\ref{tab:within}). For EARFCN~100
it moves from $-0.41$ to $0.65$ and for EARFCN~38400 from $0.67$ to $0.88$. A
model that does not know the band can only learn one averaged mapping and is
systematically offset in every band; once the band is known, the band level and
the within-band variation can be learned separately.

\begin{table}
\caption{Within-band $R^2$ under ST-PB-CV with and without the EARFCN feature.}
\label{tab:within}
\setlength{\tabcolsep}{3pt}
\begin{tabular}{|l|c|c|c|c|}
\hline
 & 100 & 1850 & 1300 & 38400\\
\hline
RSRP only & $-0.41$ & 0.65 & 0.40 & 0.67\\
$+$ EARFCN & 0.65 & 0.74 & 0.57 & 0.88\\
\hline
\end{tabular}
\end{table}

\Figure[t!](topskip=0pt, botskip=0pt, midskip=0pt){figures/fig_residual_band.pdf}
{\textbf{Residual distribution per band on the full random-hold-out test set
(9{,}262 samples).}
Boxes span the interquartile range, whiskers $1.5\times$IQR, diamonds the mean
(bias), and the dashed line is zero. The band-agnostic model is displaced in
every band; adding EARFCN centres all four distributions on zero and narrows
them.\label{fig:resid}}

Fig.~\ref{fig:resid} makes the offset explicit on all 9{,}262 samples of the
random hold-out. The
band-agnostic model is displaced in every band---by $+0.49$~dB for EARFCN~100 and
$-0.51$~dB for EARFCN~1300---and its interquartile range is wider. With EARFCN the
bias falls below $0.01$~dB in all four bands and the boxes narrow, so the feature
removes a band-dependent offset rather than merely sharpening the mapping.

We note that part of what the model learns from EARFCN is the label definition
itself (third quartile versus maximum). This is legitimate for deployment, because
the definition is fixed, but it should not be read as a purely physical gain.
To quantify the artefact we retrained under three label definitions: the mixed
definition used above (q3 for EARFCN~100, max for the others), a uniform-max
definition, and a uniform-q3 definition. These runs use a regenerated feature
table of 47{,}318 samples rather than the 46{,}306 released samples of
Tables~\ref{tab:band} and~\ref{tab:within}; their absolute MAE is therefore not
directly comparable with those tables, and only the comparison within a single
table is meaningful. Under ST-PB-CV the mixed definition gives an EARFCN gain of
22.6\% (1.07 to 0.83~dB), the uniform-max definition 15.4\% (1.00 to 0.85~dB),
and the uniform-q3 definition 31.5\% (1.13 to 0.77~dB). The mixed definition
therefore inflates the apparent feature gain by about 50\% relative to the
uniform-max baseline. More importantly, the band with the most negative
within-band $R^2$ changes with the label definition, which confirms that the
``weakest band'' is a property of the label choice rather than of the
propagation physics.

\subsection{Model Choice}
To place the protocol effect in context we compared six regressors under an
identical tuning budget (three-fold grid search on the training split) and ten
random seeds, including random forests \cite{breiman,breiman1996,liaw2002}, XGBoost \cite{chen2016} and LightGBM \cite{ke2017}, see Table~\ref{tab:models}. The spread between the best and the
worst tree ensemble is 0.15~dB, and between the best ensemble and the published
GBDT configuration 0.22~dB---well below the 0.39~dB protocol effect.

\begin{table}
\caption{Model comparison under random hold-out (ten seeds, MAE in dB). Fit
time is a single fit on one CPU core and is indicative only; it is not
normalised across machines.}
\label{tab:models}
\setlength{\tabcolsep}{3pt}
\begin{tabular}{|l|c|c|c|}
\hline
Regressor & Test MAE & $R^2$ & Fit time (s)\\
\hline
Extra trees (300) & 0.750 & 0.953 & 5\\
Random forest (300) & 0.797 & 0.950 & 15\\
XGBoost (800) & 0.852 & 0.945 & 5\\
LightGBM (800) & 0.867 & 0.944 & 2\\
GBDT, retuned & 0.902 & 0.941 & 105\\
GBDT, published \cite{zeng2024} & 0.966 & 0.934 & 52\\
Ridge regression & 1.366 & 0.882 & 0.02\\
\hline
\end{tabular}
\end{table}

\subsection{Ranking of Factors}
Combining Sections~\ref{sec:results}-A, -C and -D, the accuracy of this class of
models is governed by, in decreasing order of influence, the evaluation protocol
(0.39~dB), the availability of band information (0.32~dB) and the choice of model
(0.15--0.22~dB). Reporting a model without stating the protocol is therefore
uninformative.

\section{Device Calibration}
\label{sec:device}
Crowdsourcing implies heterogeneous handsets. RSRP carries a device-dependent
offset of several dB arising from antenna gain, front-end loss and firmware
reporting conventions, and the model has only ever seen the reference handset.

\subsection{Inter-Device Offset}
To estimate the impact we add a constant offset $\delta$ to the RSRP features of
the test set (positional features shift, dispersion features scale, and the
coefficient of variation is recomputed) and evaluate under ST-PB-CV.
Table~\ref{tab:dev} shows that a 2~dB offset increases MAE by 30\% and a 5~dB
offset by 148\%. A 100-sample output-bias calibration, estimated per band from
labelled samples disjoint from the test folds, recovers 65--75\% of the loss.

\begin{table}
\caption{Inter-device RSRP offset $\delta$ (ST-PB-CV, five folds). MAE in dB.}
\label{tab:dev}
\setlength{\tabcolsep}{3pt}
\begin{tabular}{|l|c|c|c|c|c|}
\hline
$\delta$ (dB) & 0 & 1 & 2 & 3 & 5\\
\hline
Uncalibrated & 1.046 & 1.133 & 1.360 & 1.718 & 2.585\\
Calibrated (100 samples) & 1.048 & 1.079 & 1.149 & 1.256 & 1.431\\
\hline
\end{tabular}
\end{table}

\Figure[t!](topskip=0pt, botskip=0pt, midskip=0pt){figures/fig4_device.pdf}
{\textbf{Cost of an inter-device offset and benefit of a 100-sample output
calibration.}\label{fig:dev}}

\subsection{Why Same-Device Input Calibration Does Not Work}
A natural proposal is to characterise the radio-frequency response of a given
handset---antenna gain versus frequency, front-end linearity, analogue-to-digital
conversion bias---and to correct RSRP before inversion. We tested this directly:
applying a fixed affine correction to all RSRP features, consistently to training
and test data, changes MAE from 1.0464 to 1.0458~dB, a difference of
$-0.001$~dB. The reason is structural: a tree ensemble compares feature values
only, so it is invariant to any monotone per-feature transformation, and the
correction is absorbed by the model.

Calibration therefore pays off only across devices, where the offset is a genuine
distribution shift that the model has never seen, and only for the part of the
response that is \emph{not} a monotone transformation---for example front-end
compression at high or low levels, or noise that a correction can actually remove.

\subsection{Evaluation Discipline}
A calibration procedure must not be fitted on the data used for evaluation. We
quantified the effect: estimating the per-band output bias from the test fold
itself yields 0.972~dB, whereas estimating it from disjoint labelled samples
yields 0.984~dB. The 0.012~dB difference is the leakage in this particular
setting, and it grows with the number of calibration parameters. We recommend
either a separate calibration split, or nesting the calibration inside each
training fold.

\section{Discussion}
\label{sec:discussion}

\subsection{What the Numbers Mean for Crowdsourcing}
The three protocols answer three different deployment questions. A random
hold-out answers ``can the model interpolate when neighbours are known''; ST-PB-CV
answers ``can it predict at a location or time that was never measured''; LOGO-CV
answers ``can it be used in a new campaign without recalibration''. For a
crowdsourcing platform the second and third are the relevant ones, and their MAE
is 1.05 and 1.35~dB respectively, not 0.97~dB.

\subsection{Recommendations for Reporting}
Based on the results above we suggest that future work in this area (i) reports at
least one blocked protocol in addition to the random split, (ii) states the
feature set explicitly, in particular whether the band identity is used, (iii)
reports MAE and RMSE together with a pooled $R^2$, and (iv) states the handset
model and whether any cross-device calibration was applied.

\subsection{Limitations}
All measurements come from a single handset and a single campus area of
$380\times260$~m, so the cross-device results of Section~\ref{sec:device} are
obtained by simulation rather than by measurement, and the conclusions may not
transfer to indoor or dense-urban environments. The dataset contains only RSRP,
so we cannot test whether the gain obtained from EARFCN is complementary to, or
overlapping with, the gain reported for multi-parameter inputs \cite{huang2026}.
Finally, the block size and embargo width were selected on physical grounds.
A brief sensitivity analysis shows that widening the embargo from
0~m/0~s to 20~m/60~s and then to 50~m/300~s increases MAE from
1.03 to 1.05 and 1.07~dB, confirming that the residual leakage
across the 20~m/60~s boundary is small.

\section{Conclusion}
\label{sec:conclusion}
On a public dataset we reproduced a representative smartphone-based EMF inversion
model and then varied three factors that are usually fixed by convention. The
evaluation protocol is the most influential: random splitting understates the
error by 8\% relative to spatio-temporal blocking and by 40\% relative to holding
out a whole campaign. Band information, which the handset already provides, is the
cheapest genuine improvement, worth 15--31\% depending on how the label is
defined. Device calibration matters only across
devices, and a few hundred labelled samples suffice. Under the strictest protocol
the band-aware model still outperforms the originally reported random-split
accuracy by 10\%. We therefore suggest that protocol, feature set and device
identity be reported together in future work on crowdsourced EMF monitoring.

\appendices
\section{Recovering Unreported Implementation Details}
\label{app:details}
Three details absent from \cite{zeng2024} had to be recovered from the released
sample table.
\begin{enumerate}
\item \emph{EWMA coefficient.} The paper states only that $\alpha$ lies between 0
and 1. Fitting the released features gives $\alpha=0.05$; values in the range
0.03--0.10 reproduce the released statistics to within 1~dB.
\item \emph{Quantile definition.} The quartiles are order statistics at positions
$(N+1)/4$, $(N+1)/2$ and $3(N+1)/4$ of the seven-point window, not linearly
interpolated quantiles.
\item \emph{Window composition.} The seven appended raw values are unfiltered
linear interpolations on the 1~s grid, whereas the 16 statistics are computed on
the EWMA-filtered series.
\end{enumerate}
With these choices the reconstructed feature table matches the released one
exactly for the seven raw values and with a correlation of 0.988 for the
positional statistics.

\section*{Data and Code Availability}
The dataset analysed in this study is publicly available from figshare at
\url{https://doi.org/10.6084/m9.figshare.25034537.v1}~\cite{zeng2024}. The
feature-engineering, cross-validation and calibration code, together with the
scripts that regenerate every table and figure of this paper, are openly
available at \url{https://doi.org/10.5281/zenodo.23758133} (development version:
\url{https://github.com/347661070zlz-cloud/emf-inversion-audit}).

\section*{Acknowledgment}
The authors thank the authors of \cite{zeng2024} for releasing their dataset and
code, which made this study possible.

\begin{thebibliography}{22}

\bibitem{aerts2013}
S.~Aerts, D.~Deschrijver, L.~Verloock, T.~Dhaene, L.~Martens, and W.~Joseph,
``Assessment of outdoor radiofrequency electromagnetic field exposure through
hotspot localization using kriging-based sequential sampling,''
\emph{Environmental Research}, vol.~126, pp.~184--191, 2013,
doi:~10.1016/j.envres.2013.05.005.

\bibitem{ganti2011}
R.~K. Ganti, F.~Ye, and H.~Lei, ``Mobile crowdsensing: Current state and
future challenges,'' \emph{IEEE Commun. Mag.}, vol.~49, no.~11,
pp.~32--39, Nov. 2011, doi:~10.1109/MCOM.2011.6069707.

\bibitem{lane2010}
N.~D. Lane, E.~Miluzzo, H.~Lu, D.~Peebles, T.~Choudhury, and A.~T. Campbell,
``A survey of mobile phone sensing,'' \emph{IEEE Commun. Mag.}, vol.~48,
no.~9, pp.~140--150, Sep. 2010, doi:~10.1109/MCOM.2010.5560598.

\bibitem{icnirp}
International Commission on Non-Ionizing Radiation Protection, ``Guidelines for
limiting exposure to electromagnetic fields (100~kHz to 300~GHz),''
\emph{Health Physics}, vol.~118, no.~5, pp.~483--524, 2020,
doi:~10.1097/HP.0000000000001210.

\bibitem{ieee2019}
IEEE, ``IEEE Standard for Safety Levels with Respect to Human Exposure to
Electric, Magnetic, and Electromagnetic Fields, 0~Hz to 300~GHz,''
IEEE Std C95.1-2019, 2019, doi:~10.1109/IEEESTD.2019.8859679.

\bibitem{schie}
S.~Schie{\ss}l, T.~Ta, T.~Kopacz, and D.~Heberling, ``Konzept und Messungen
f\"ur ein Smartphone-basiertes Immissions-Monitoring-System,'' in \emph{Proc.
EMV: Int. Fachmesse und Kongress f\"ur Elektromagnetische
Vertr\"aglichkeit}, Cologne, Germany, 2020, doi:~10.15488/10060.

\bibitem{zeng2024}
S.~Zeng, W.~Chen, Y.~Ji, L.~Yan, and X.~Zhao, ``Measurement and calibration of
EMF: A study using phone and GBDT for mobile communication signals,''
\emph{Radio Science}, vol.~59, no.~2, p.~e2023RS007890, 2024,
doi:~10.1029/2023RS007890.

\bibitem{friedman}
J.~H. Friedman, ``Greedy function approximation: A gradient boosting machine,''
\emph{Ann. Statist.}, vol.~29, no.~5, pp.~1189--1232, 2001,
doi:~10.1214/aos/1013203451.

\bibitem{friedman2002}
J.~H. Friedman, ``Stochastic gradient boosting,'' \emph{Comput. Statist.
Data Anal.}, vol.~38, no.~4, pp.~367--378, 2002,
doi:~10.1016/S0167-9473(01)00065-2.

\bibitem{hastie2009}
T.~Hastie, R.~Tibshirani, and J.~Friedman, \emph{The Elements of Statistical
Learning: Data Mining, Inference, and Prediction}, 2nd~ed.\ New York, NY, USA:
Springer, 2009, doi:~10.1007/978-0-387-84858-7.

\bibitem{huang2026}
Q.~Huang, S.~Zeng, X.~Zhao, and L.~Yan, ``Research on measurement and perception
of wireless communication electromagnetic environment based on multiple smartphone
parameters,'' \emph{Chinese J. Radio Sci.}, 2026,
doi:~10.12265/j.cjors.2026025.

\bibitem{roberts}
D.~R. Roberts, V.~Bahn, S.~Ciuti, M.~S. Boyce, J.~Elith, G.~Guillera-Arroita,
\emph{et al.}, ``Cross-validation strategies for data with temporal, spatial,
hierarchical, or phylogenetic structure,'' \emph{Ecography}, vol.~40, no.~8,
pp.~913--929, 2017, doi:~10.1111/ecog.02881.

\bibitem{kara}
N.~Karasiak, J.-F. Dejoux, C.~Monteil, and D.~Sheeren, ``Spatial dependence
between training and test sets: Another pitfall of classification accuracy
assessment in remote sensing,'' \emph{Machine Learning}, vol.~111, no.~7,
pp.~2715--2740, 2022, doi:~10.1007/s10994-021-05972-1.

\bibitem{pohjankukka2017}
J.~Pohjankukka, T.~Pahikkala, P.~Nevalainen, and J.~Heikkonen, ``Estimating the
prediction performance of spatial models via spatial $k$-fold cross
validation,'' \emph{Int. J. Geographical Inf. Sci.}, vol.~31, no.~10,
pp.~2001--2019, 2017, doi:~10.1080/13658816.2017.1346255.

\bibitem{meyer2021}
H.~Meyer and E.~Pebesma, ``Predicting into unknown space? Estimating the area
of applicability of spatial prediction models,'' \emph{Methods Ecol. Evol.},
vol.~12, no.~9, pp.~1620--1633, 2021, doi:~10.1111/2041-210X.13650.

\bibitem{valavi2019}
R.~Valavi, J.~Elith, J.~J. Lahoz-Monfort, and G.~Guillera-Arroita, ``blockCV: An
R package for generating spatially or environmentally separated folds for
$k$-fold cross-validation of species distribution models,''
\emph{Methods Ecol. Evol.}, vol.~10, no.~2, pp.~225--232, 2019,
doi:~10.1111/2041-210X.13107.

\bibitem{3gpp}
3GPP, ``Evolved Universal Terrestrial Radio Access (E-UTRA); Physical
layer; Measurements,'' 3GPP TS 36.214, V17.2.0, 2022.

\bibitem{breiman}
L.~Breiman, ``Random forests,'' \emph{Machine Learning}, vol.~45, no.~1,
pp.~5--32, 2001, doi:~10.1023/A:1010933404324.

\bibitem{breiman1996}
L.~Breiman, ``Bagging predictors,'' \emph{Machine Learning}, vol.~24,
no.~2, pp.~123--140, 1996, doi:~10.1007/BF00058655.

\bibitem{liaw2002}
A.~Liaw and M.~Wiener, ``Classification and regression by randomForest,''
\emph{R News}, vol.~2, no.~3, pp.~18--22, 2002.

\bibitem{chen2016}
T.~Chen and C.~Guestrin, ``XGBoost: A scalable tree boosting system,''
in \emph{Proc. 22nd ACM SIGKDD Int. Conf. Knowl. Discovery Data Mining},
2016, pp.~785--794, doi:~10.1145/2939672.2939785.

\bibitem{ke2017}
G.~Ke, Q.~Meng, T.~Finley, T.~Wang, W.~Chen, W.~Ma, Q.~Ye, and T.-Y. Liu,
``LightGBM: A highly efficient gradient boosting decision tree,''
in \emph{Adv. Neural Inf. Process. Syst.}, vol.~30, 2017, pp.~3146--3154.

\end{thebibliography}

\begin{IEEEbiographynophoto}{Juan Pan}
is currently a Professor-level Senior Engineer and
the Deputy Director of the CTTL System Laboratory, China Academy of Information
and Communications Technology, Beijing, China. Her research interests include
electromagnetic compatibility and radio-frequency exposure assessment.
\end{IEEEbiographynophoto}

\begin{IEEEbiographynophoto}{Lianzhong Zheng}
is currently pursuing the M.S. degree with
the China Academy of Information and Communications Technology, Beijing, China.
His research interests include electromagnetic compatibility and machine learning
for radio-frequency exposure assessment.
\end{IEEEbiographynophoto}

\EOD

\end{document}
